% 11-ownership.tex — Clause 11: Ownership and linearity
\chapter{Ownership and linearity}
\label{sec:11-ownership}
\index{ownership}
\index{linearity}

\pnum
This clause specifies Lain's linear type system, ownership modes, borrow
checker, and non-lexical lifetime rules.  All rules in this clause are
enforced at compile time with zero runtime overhead.

\section{Ownership modes}
\label{sec:own.modes}
\index{ownership mode}

\pnum
Every parameter, return value, and struct field has one of three ownership
modes:

\begin{longtable}{@{}llll@{}}
\toprule
\textbf{Mode} & \textbf{Param syntax} & \textbf{Call-site} & \textbf{Semantics} \\
\midrule
\endfirsthead
\toprule
\textbf{Mode} & \textbf{Param syntax} & \textbf{Call-site} & \textbf{Semantics} \\
\midrule
\endhead
\bottomrule
\endlastfoot
Shared  & \lain{name T}     & \lain{f(x)}     & read-only; multiple concurrent borrows allowed \\
Mutable & \lain{var name T} & \lain{f(var x)} & exclusive read-write; no other borrow concurrent \\
Owned   & \lain{mov name T} & \lain{f(mov x)} & transfers ownership; source becomes Consumed \\
\end{longtable}

\pnum
Local variables declared with \lain{var name = expr} have mutable storage
accessible within their scope; their ownership mode in the type system is
Owned (they hold linear values that must be consumed at scope exit).

\section{Linear states}
\label{sec:own.states}
\index{linear state}

\pnum
Every variable of a linear type has a \textit{linear state} at each program
point.  The three states are:

\begin{description}[leftmargin=2em]
  \item[Uninit]
    The variable has been declared but no value has been assigned.
  \item[Unconsumed]
    The variable holds a value that has not yet been transferred.
  \item[Consumed]
    The variable's value has been transferred via \lain{mov} or a consuming
    function call.
\end{description}

\pnum
The state machine governing transitions is:

\begin{center}
\begin{tabular}{lll}
\toprule
\textbf{State before} & \textbf{Operation} & \textbf{State after} \\
\midrule
Uninit      & first assignment              & Unconsumed \\
Unconsumed  & \lain{mov x} (transfer)       & Consumed \\
Unconsumed  & borrow (\lain{var} or shared) & Unconsumed (borrow active) \\
Consumed    & assignment (overwrite)        & Unconsumed \\
Consumed    & read or \lain{mov}            & \illformed{} \\
Uninit      & read or \lain{mov}            & \illformed{} \\
\bottomrule
\end{tabular}
\end{center}

\begin{note}
  The Consumed~$\to$~Unconsumed transition via overwrite enables the
  loop-reassign-then-consume pattern: a linear variable declared outside a
  loop may be consumed and reassigned on each iteration without leaking.
\end{note}

\section{Diagnostics}
\label{sec:own.diag}
\index{ownership diagnostics}

\begin{constraint}
  Using a variable that is in the Consumed or Uninit state (reading its
  value or transferring it with \lain{mov}) is \illformed{}~(\diagcode{E001}).
\end{constraint}

\begin{constraint}
  Transferring ownership of a variable that has already been transferred
  (double-move) is \illformed{}~(\diagcode{E002}).
\end{constraint}

\begin{constraint}
  A linear variable that is in the Unconsumed state when its scope exits
  (resource leak) is \illformed{}~(\diagcode{E003}).
\end{constraint}

\begin{constraint}
  An aliasing violation — attempting a mutable borrow when a shared borrow
  is active, or a shared borrow when a mutable borrow is active, or a
  second mutable borrow of the same variable — is
  \illformed{}~(\diagcode{E004}).
\end{constraint}

\begin{constraint}
  Attempting to move a variable that has an active borrow is
  \illformed{}~(\diagcode{E005}).
\end{constraint}

\begin{constraint}
  Moving a variable inside a loop body where the loop may execute more than
  once is \illformed{}~(\diagcode{E006}).
\end{constraint}

\begin{constraint}
  If a linear variable is consumed in one branch of an \lain{if}/\lain{else}
  chain but not in all branches, the program is
  \illformed{}~(\diagcode{E007}).
\end{constraint}

\begin{constraint}
  Moving a value that is currently borrowed (\lain{var} or shared), or
  moving a struct after some of its fields have already been consumed,
  is \illformed{}~(\diagcode{E008}).
\end{constraint}

\begin{constraint}
  If a linear value is in state \textsc{Consumed} on some control-flow
  paths through an \lain{if}/\lain{case} statement and in state
  \textsc{Unconsumed} on others, the merged state at the join point
  is undefined and the program is \illformed{}~(\diagcode{E016}). The
  same diagnostic applies at field level: if a linear field of a
  struct is consumed on some paths but not on others, \diagcode{E016}
  is issued for that field.
\end{constraint}

\section{Mutable parameters for primitive types (\texttt{var T})}
\label{sec:own.var-primitive}
\index{var parameter!primitive type}
\index{write-through}

\pnum
A parameter declared \lain{var name T}, where \lain{T} is a primitive type
(\lain{iN}, \lain{uN}, \lain{f32}, \lain{f64}, \lain{bool}), is passed
\textbf{by pointer} in the C ABI: the parameter's C type is \texttt{T*}.

\pnum
\textbf{Write-through semantics}: an assignment to such a parameter inside
the function body writes through the pointer and is visible to the caller.
Reading the parameter's value automatically dereferences the pointer.

\begin{example}
\begin{laincode}
proc increment(var n i32) {
    n = n + 1    // writes through: caller's variable is updated
}

proc main() i32 {
    var x = 0
    increment(var x)
    return x     // x == 1
}
\end{laincode}
\end{example}

\pnum
\textbf{Forwarding}: when a \lain{var T} primitive parameter is passed as
argument to another \lain{var T} parameter, the pointer is forwarded
directly without dereferencing and re-taking the address.

\pnum
This behavior is identical to mutable parameters of struct type, where the
by-pointer ABI was already in effect.  The distinction is that for primitive
types the codegen previously did not dereference reads; this was a bug that
has been corrected.

\begin{note}
  The call site shall annotate the argument with \lain{var}:
  \lain{f(var x)}.  Omitting \lain{var} passes a shared (read-only)
  borrow, which does not permit writes through the callee.
\end{note}

\section{The read-write lock invariant}
\label{sec:own.rwlock}
\index{read-write lock invariant}

\pnum
\textbf{Invariant}: At every program point, for every variable~$V$, exactly
one of the following holds:
\begin{enumerate}
  \item $V$ has $N \geq 0$ active shared borrows and zero active mutable borrows.
  \item $V$ has exactly one active mutable borrow and zero shared borrows.
\end{enumerate}

\pnum
This invariant is the compile-time analogue of a reader-writer lock.  A
conforming implementation shall enforce it at every program point.

\section{Non-lexical lifetimes (NLL)}
\label{sec:own.nll}
\index{non-lexical lifetime}

\pnum
A borrow's lifetime ends at the borrow's \textit{last use}, not at the end
of the variable's lexical scope.  The implementation computes last-use
points via a use-analysis pre-pass over the AST.

\begin{example}
\begin{laincode}
var data = Data(42)
var r = get_ref(var data)  // mutable borrow starts
use_ref(var r)             // last use of r: borrow ends HERE (NLL)
var y = read(data)         // OK: borrow has expired
\end{laincode}
\end{example}

\section{Two-phase borrows}
\label{sec:own.two-phase}
\index{two-phase borrow}

\pnum
During argument evaluation for a function call, borrows of the receiver
object are placed in the \textit{RESERVED} phase.  A RESERVED borrow
allows concurrent shared reads of the same owner but blocks mutable writes
and moves.  After all arguments have been evaluated, RESERVED borrows are
promoted to the \textit{ACTIVE} phase.

\pnum
This mechanism enables UFCS calls of the form \lain{x.f(x.field)} where
the first argument (the receiver) requires a mutable borrow of \lain{x}
while the second argument reads a field of \lain{x}.

\begin{longtable}{@{}lll@{}}
\toprule
\textbf{Phase} & \textbf{Allows on owner} & \textbf{Blocks on owner} \\
\midrule
\endfirsthead
\toprule
\textbf{Phase} & \textbf{Allows on owner} & \textbf{Blocks on owner} \\
\midrule
\endhead
\bottomrule
\endlastfoot
RESERVED & shared reads & mutable writes, moves \\
ACTIVE   & (nothing)    & all other access \\
\end{longtable}

\begin{example}
\begin{laincode}
proc push_n(var v Vec, n int) { v.data = v.data + n }

var v = Vec(0, 10)
v.push_n(v.cap)   // desugars to push_n(var v, v.cap)
                  // var v enters RESERVED phase; v.cap read is allowed
\end{laincode}
\end{example}

\section{Move semantics}
\label{sec:own.move}
\index{move semantics}

\pnum
The \lain{mov} operator transfers ownership.  The source variable
transitions to Consumed state; any subsequent read or move of the source
is \illformed{}~(\diagcode{E001}).

\pnum
Ownership transfer is the only way to pass a linear value to a consuming
parameter.  The \lain{mov} keyword shall be present at the call site:
\lain{f(mov x)}.

\section{Borrowed match}
\label{sec:own.borrowed-match}
\index{borrowed match}

\pnum
The \lain{case \&expr \{ ... \}} form (see \S\,\ref{sec:15-match}) performs
pattern matching without consuming the scrutinee.  A temporary shared
borrow of the scrutinee is registered for the duration of the match body
and released when the match completes.

\begin{constraint}
  Inside a \lain{case \&expr} body, the scrutinee shall not be mutated
  or moved.  Attempting to do so violates the active shared
  borrow~(\diagcode{E004}).
\end{constraint}

\section{Loop and ownership}
\label{sec:own.loop}
\index{ownership in loops}

\pnum
A linear variable declared outside a loop may be consumed inside the loop
body if and only if it is also reassigned (overwritten) on every path
through the loop body before the next iteration begins.

\begin{example}
\begin{laincode}
import std.fs

proc main() int {
    var f = open_file("data.txt", "r")
    close_file(mov f)        // consumed before loop
    var i = 0
    while i < 3 decreasing 3 - i {
        f = open_file("data.txt", "r")  // reassign (Consumed -> Unconsumed)
        close_file(mov f)               // consume
        i = i + 1
    }
    return 0
}
\end{laincode}
\end{example}

\section{Branch consistency}
\label{sec:own.branches}
\index{branch consistency}

\pnum
After an \lain{if}/\lain{else} chain or \lain{case} match, the linear
state of each variable from the enclosing scope is the
\textit{intersection} of its states across all branches:

\begin{itemize}
  \item If consumed in all branches: the variable is Consumed after the chain.
  \item If consumed in no branches: the variable is Unconsumed after the chain.
  \item If consumed in some but not all branches: \illformed{}~(\diagcode{E016}).
\end{itemize}

\pnum
The same rule applies at field level: if a linear field of a struct is
consumed on some branches but not on others, \diagcode{E016} is emitted
for that field.

\begin{example}
\begin{laincode}
type Handle { mov ptr *i32 }

func maybe_consume(mov h Handle, flag bool) {
    if flag {
        drop(mov h)   // consumed in then-branch
    }
    // [E016]: h consumed in then but not in else (implicit else)
}
\end{laincode}
\end{example}

\begin{note}
  \diagcode{E007} (generic ownership mismatch) was the diagnostic emitted
  before \diagcode{E016} was introduced.  \diagcode{E016} is now the
  dedicated code for branch-inconsistency on linear values.
\end{note}

\section{Field-level tracking}
\label{sec:own.fields}
\index{field-level borrow tracking}

\pnum
The borrow checker tracks borrows at the field level.  Borrowing
non-overlapping fields of the same struct simultaneously is permitted:

\begin{laincode}
type Pair { a i32, b i32 }
var p = Pair(1, 2)
var ra = read_a(p.a)   // shared borrow of p.a (implicit)
var rb = read_b(p.b)   // shared borrow of p.b -- OK: non-overlapping
\end{laincode}

\begin{note}
  The reference implementation tracks field borrows at one level of depth.
  Nested field borrows (e.g., \lain{p.a.x} vs \lain{p.a.y}) are conservatively
  treated as overlapping.
\end{note}
